% GENERATED FILE. Edit canonical lessons, then run npm run generate:books.
% canonical-source-hash: fnv1a64:bce8001e29557a2a
% canonical-lessons: HI-C134-botal, HI-C134-pensil, HI-C134-tasvir, HI-C134-sisa, HI-C134-anguthi

\chapter{Bottle, Pencil, Picture, Mirror, Ring}
\label{ch:hi-bottle-pencil-picture-mirror-ring}
\hlchaptermodality{134}

\begin{chapteropening}
\textbf{By the end of this chapter:} \emph{I can say a bottle, a pencil, a picture, a mirror and a ring.}

Complete the last lesson of chapter 134: i can say a bottle, a pencil, a picture, a mirror and a ring.
\end{chapteropening}

\section[botal]{\hi{बोतल} (botal) — a bottle}

\label{lesson:HI-C134-botal}



\begin{quote}
Before the new one: say the Hindi for a bed frame, then the Hindi for bedding.

Order three things two ways: \textbf{\hi{पहली}}, \textbf{\hi{दूसरी}}, \textbf{\hi{तीसरी}}, then \textbf{\hi{एक} … \hi{दूसरा}}. Count on to \textbf{\hi{नौवाँ}} and \textbf{\hi{दसवाँ}}.
\end{quote}

\subsection*{You'll want to know: \hi{बोतल}}
\textbf{\hi{बोतल}} — \emph{botal} — \textquotedblleft{}a bottle\textquotedblright{}.

\begin{grammarlens}[title={What you've built}]
One more word for this chapter.
\end{grammarlens}

\subsection*{Your turn}
\begin{itemize}
\raggedright
  \item \emph{Say it:} \emph{botal}
  \item \emph{Say it:} \emph{botal}, once more
  \item \emph{Say it:} say \emph{botal}
  \item \emph{Recall:} say the Hindi for a bed frame, then the Hindi for bedding, then say \emph{botal} again
\end{itemize}

\begin{tcolorbox}[breakable,colback=teal!4,colframe=teal!35!black,title={Before you move on}]
What does \hi{बोतल} mean? (\textquotedblleft{}A bottle\textquotedblright{}.) Say it once more.
\end{tcolorbox}
\section[pensil]{\hi{पेंसिल} (pensil) — a pencil}

\label{lesson:HI-C134-pensil}



\begin{quote}
Before the new one: say the Hindi for bedding, then the Hindi for a bottle.

Say the four mind verbs and what each first meant: \textbf{\hi{लिखना}}, to scratch; \textbf{\hi{पढ़ना}}, to recite; \textbf{\hi{समझना}}, to wake up to; \textbf{\hi{सोचना}}, to burn.
\end{quote}

\subsection*{You'll want to know: \hi{पेंसिल}}
\textbf{\hi{पेंसिल}} — \emph{pensil} — \textquotedblleft{}a pencil\textquotedblright{}.

\begin{grammarlens}[title={What you've built}]
One more word for this chapter.
\end{grammarlens}

\subsection*{Your turn}
\begin{itemize}
\raggedright
  \item \emph{Say it:} \emph{pensil}
  \item \emph{Say it:} \emph{pensil}, once more
  \item \emph{Say it:} say \emph{pensil}
  \item \emph{Recall:} say the Hindi for bedding, then the Hindi for a bottle, then say \emph{pensil} again
\end{itemize}

\begin{tcolorbox}[breakable,colback=teal!4,colframe=teal!35!black,title={Before you move on}]
What does \hi{पेंसिल} mean? (\textquotedblleft{}A pencil\textquotedblright{}.) Say it once more.
\end{tcolorbox}
\section[tasvīr]{\hi{तस्वीर} (tasvīr) — a picture}

\label{lesson:HI-C134-tasvir}



\begin{quote}
Before the new one: say the Hindi for a bottle, then the Hindi for a pencil.
\end{quote}

\subsection*{You'll want to know: \hi{तस्वीर}}
\textbf{\hi{तस्वीर}} — \emph{tasvīr} — \textquotedblleft{}a picture\textquotedblright{}.

\begin{grammarlens}[title={What you've built}]
One more word for this chapter.
\end{grammarlens}

\subsection*{Your turn}
\begin{itemize}
\raggedright
  \item \emph{Say it:} \emph{tasvīr}
  \item \emph{Say it:} \emph{tasvīr}, once more
  \item \emph{Say it:} say \emph{tasvīr}
  \item \emph{Recall:} say the Hindi for a bottle, then the Hindi for a pencil, then say \emph{tasvīr} again
\end{itemize}

\begin{tcolorbox}[breakable,colback=teal!4,colframe=teal!35!black,title={Before you move on}]
What does \hi{तस्वीर} mean? (\textquotedblleft{}A picture\textquotedblright{}.) Say it once more.
\end{tcolorbox}
\section[śīśā]{\hi{शीशा} (śīśā) — a mirror, glass}

\label{lesson:HI-C134-sisa}



\begin{quote}
Before the new one: say the Hindi for a pencil, then the Hindi for a picture.
\end{quote}

\subsection*{You'll want to know: \hi{शीशा}}
\textbf{\hi{शीशा}} — \emph{śīśā} — \textquotedblleft{}a mirror, glass\textquotedblright{}.

\begin{grammarlens}[title={What you've built}]
One more word for this chapter.
\end{grammarlens}

\subsection*{Your turn}
\begin{itemize}
\raggedright
  \item \emph{Say it:} \emph{śīśā}
  \item \emph{Say it:} \emph{śīśā}, once more
  \item \emph{Say it:} say \emph{śīśā}
  \item \emph{Recall:} say the Hindi for a pencil, then the Hindi for a picture, then say \emph{śīśā} again
\end{itemize}

\begin{tcolorbox}[breakable,colback=teal!4,colframe=teal!35!black,title={Before you move on}]
What does \hi{शीशा} mean? (\textquotedblleft{}A mirror, glass\textquotedblright{}.) Say it once more.
\end{tcolorbox}
\section[aṅgūṭhī]{\hi{अंगूठी} (aṅgūṭhī) — a ring}

\label{lesson:HI-C134-anguthi}



\begin{quote}
Before the new one: say the Hindi for a picture, then the Hindi for a mirror.
\end{quote}

\subsection*{You'll want to know: \hi{अंगूठी}}
\textbf{\hi{अंगूठी}} — \emph{aṅgūṭhī} — \textquotedblleft{}a ring\textquotedblright{}.

\begin{grammarlens}[title={What you've built}]
One more word for this chapter.
\end{grammarlens}

\subsection*{Your turn}
\begin{itemize}
\raggedright
  \item \emph{Say it:} \emph{aṅgūṭhī}
  \item \emph{Say it:} \emph{aṅgūṭhī}, once more
  \item \emph{Say it:} say \emph{aṅgūṭhī}
  \item \emph{Recall:} say the Hindi for a picture, then the Hindi for a mirror, then say \emph{aṅgūṭhī} again
\end{itemize}

\begin{tcolorbox}[breakable,colback=teal!4,colframe=teal!35!black,title={Before you move on}]
What does \hi{अंगूठी} mean? (\textquotedblleft{}A ring\textquotedblright{}.) Say it once more.
\end{tcolorbox}
